% ============================================================================
\chapter{Introdução}
% ============================================================================

% Contextualização do tema
Contextualizar o tema do trabalho, apresentando o cenário geral e a
relevância do assunto.

% ============================================================================
\section{Problema de Pesquisa}
% ============================================================================

Apresentar o problema que motivou a pesquisa, formulando-o preferencialmente
como uma pergunta.

% ============================================================================
\section{Objetivos}
% ============================================================================

\subsection{Objetivo Geral}

Descrever o objetivo geral do trabalho de forma clara e concisa.

\subsection{Objetivos Específicos}

\begin{itemize}
    \item Primeiro objetivo específico;
    \item Segundo objetivo específico;
    \item Terceiro objetivo específico.
\end{itemize}

% ============================================================================
\section{Justificativa}
% ============================================================================

Justificar a importância e relevância do trabalho, destacando as
contribuições esperadas.

% ============================================================================
\section{Estrutura do Trabalho}
% ============================================================================

Este trabalho está organizado em cinco capítulos. O Capítulo~\ref{chap:fundamentacao}
apresenta a fundamentação teórica. O Capítulo~\ref{chap:metodologia} descreve
a metodologia utilizada. O Capítulo~\ref{chap:resultados} apresenta os
resultados obtidos. Por fim, o Capítulo~\ref{chap:conclusao} traz as
considerações finais.
